\documentclass{article}
\usepackage{amsmath,amssymb}
\usepackage{xurl}
\usepackage[hidelinks]{hyperref}
\setlength{\emergencystretch}{2em}
% Shared commands for notes in docs/paper-gaps/.

\DeclareUnicodeCharacter{2080}{\ifmmode{}_{0}\else\textsubscript{0}\fi}
\DeclareUnicodeCharacter{2081}{\ifmmode{}_{1}\else\textsubscript{1}\fi}
\DeclareUnicodeCharacter{2082}{\ifmmode{}_{2}\else\textsubscript{2}\fi}
\DeclareUnicodeCharacter{2099}{\ifmmode{}_{n}\else\textsubscript{n}\fi}

\newcommand{\Lean}{\textsc{Lean}}
\ExplSyntaxOn
\NewDocumentCommand{\leanid}{m}{
  \ifmmode
    \textsf{\detokenize{#1}}
  \else
    \group_begin:
    \urlstyle{sf}
    \tl_set:Nn \l_tmpa_tl {#1}
    \tl_replace_all:Nnn \l_tmpa_tl {\_} {_}
    \exp_args:NV \path \l_tmpa_tl
    \group_end:
  \fi
}
\ExplSyntaxOff
\providecommand{\tr}{\operatorname{tr}}
\newcommand{\tnleanIssueBase}{https://github.com/LionSR/TNLean/issues/}
\newcommand{\tnleanPrBase}{https://github.com/LionSR/TNLean/pull/}
\newcommand{\tnleanDocBase}{https://sirui-lu.com/TNLean/docs/}
\newcommand{\ghissue}[1]{\href{\tnleanIssueBase#1}{GitHub issue \##1}}
\newcommand{\ghpr}[1]{\href{\tnleanPrBase#1}{PR \##1}}
\newcommand{\doclink}[1]{\href{\tnleanDocBase#1.html}{\path{#1}}}

% Dirac notation and the identity operator, as in blueprint/src/macros/common.tex.
% \providecommand keeps note-local definitions authoritative.
\usepackage{dsfont}
\providecommand{\ket}[1]{|#1\rangle}
\providecommand{\bra}[1]{\langle #1|}
\providecommand{\braket}[2]{\langle #1 | #2 \rangle}
\providecommand{\ketbra}[2]{|#1\rangle\!\langle #2|}
\providecommand{\Id}{\mathds{1}}
\providecommand{\one}{\mathds{1}}
\providecommand{\C}{\mathbb{C}}
\providecommand{\MN}[1]{M_{#1}(\C)}
\providecommand{\MD}{M_D(\C)}

% External referencing for paper sources.  The build helper writes sanitized
% auxiliary files under build/paper-aux/ and excludes duplicated source labels.
\usepackage{xr}

% Import a generated paper auxiliary file with a caller-supplied label prefix.
% The candidate paths support builds from docs/paper-gaps/, the repository root,
% and build/paper-aux/ itself.
\newcommand{\importpaperaux}[2]{%
  \IfFileExists{../../build/paper-aux/#2.aux}{%
    \externaldocument[#1]{../../build/paper-aux/#2}%
  }{%
    \IfFileExists{build/paper-aux/#2.aux}{%
      \externaldocument[#1]{build/paper-aux/#2}%
    }{%
      \IfFileExists{#2.aux}{%
        \externaldocument[#1]{#2}%
      }{}%
    }%
  }%
}

% General external reference macros
\newcommand{\paperref}[2]{\ref{#1-#2}}
\newcommand{\papereqref}[2]{(\ref{#1-#2})}

% CPSV16 source (1606.00608 / MPDO-22-12-17-2.tex) external document and shortcuts
\importpaperaux{cpsv16-}{1606.00608/MPDO-22-12-17-2-tnlean}
\newcommand{\cpsvref}[1]{\paperref{cpsv16}{#1}}
\newcommand{\cpsveqref}[1]{\papereqref{cpsv16}{#1}}

% Verdict marker: \gapnote{<kind>}{<status>}. Kind is the policy
% classification (clarification, local-correction, scope-restriction,
% unfaithful, false-source, open-gap); status is open, wip (elimination
% underway), resolved, or historical. Machine-read by the site index;
% prints a verdict line.
\newcommand{\gapnote}[2]{\par\noindent\textbf{Verdict:} #1 (#2).\par}

\title{The $\mathrm{Rep}(S_3)$ and $su(2)_4$ Phases: Completeness of the Lists}
\author{}
\date{2026-09-27}
\begin{document}
\maketitle
\gapnote{scope-restriction}{open}

\section{The source claims}
Garre-Rubio, Lootens and Moln\'ar list the phases of matrix product states
symmetric under matrix product operators representing $su(2)_4$ and
$\mathrm{Rep}(S_3)$~\cite{GarreRubio2023MPOSym}. For $su(2)_4$, with labels
$0,\tfrac12,1,\tfrac32,2$, they state that there are only two phases, the one
with five blocks acted on by the fusion rules and the module
$\mathcal M_{TY}$ with four blocks $x,y,z,s$, and they print the action of
every label on the blocks of $\mathcal M_{TY}$.\footnote{Source traceability:
    \path{Papers/2203.12563/REsubmission.tex}, lines 1890--1931.} For
$\mathrm{Rep}(S_3)$, with labels $1,\pi,\psi$, they state that every phase
is characterized by a subgroup of $S_3$, and print the actions on the blocks
for $H=\mathbb Z_1,\mathbb Z_2,\mathbb Z_3,S_3$. They identify the
$\mathbb Z_1$, $\mathbb Z_2$ and $\mathbb Z_3$ phases with restrictions of
the $su(2)_4$ phases along the embedding $1,\pi,\psi\mapsto0,1,2$: the block
$s$ of $\mathcal M_{TY}$, the blocks $\tfrac12,\tfrac32$ of the first phase,
and the blocks $x,y,z$ of $\mathcal M_{TY}$.\footnote{Source traceability:
    lines 1933--1989; the restrictions at lines 1927, 1940, 1957 and 1972.}
The multiplicities $M_{a,x}^y$ of each phase form a nonnegative integer
representation of the fusion ring,
\begin{align}
    \sum_c N_{ab}^c M_{c,x}^y & =\sum_z M_{b,x}^z M_{a,z}^y ,
    \label{eq:glm23_list_nimrep}
\end{align}
with the unit acting as the identity.

\section{Completeness of the lists}
Formalized: every listed action satisfies (\ref{eq:glm23_list_nimrep}), with the $\mathbb Z_3$ action
corrected as recorded in the note on the $\mathbb Z_3$ table, the restrictions
stated by the source hold, and on a single block the list is complete. A representation on one
block with the unit acting trivially is a ring homomorphism to $\mathbb Z$
that is nonnegative on the labels; for $\mathrm{Rep}(S_3)$ the only one is
$(1,2,1)$, the $\mathbb Z_1$ phase, and $su(2)_4$ has none, since it would force
$m_2=1$, $m_1=2$ and $m_{1/2}^2=3$.\footnote{Results:
    \leanid{MPOTensor.isFusionCharacter_repS3Fusion_iff} and
    \leanid{MPOTensor.not_isFusionCharacter_su24Fusion}.}

Not formalized: completeness of the lists beyond one block. The source takes
it from the classification of module categories, which it cites. It is not a
consequence of (\ref{eq:glm23_list_nimrep}) with trivial unit action alone:
for $\mathrm{Rep}(S_3)$ on two blocks,
\begin{align}
    M_\psi & =\begin{pmatrix}1&0\\0&1\end{pmatrix}, &
    M_\pi  & =\begin{pmatrix}0&1\\2&1\end{pmatrix}
    \notag
\end{align}
satisfies (\ref{eq:glm23_list_nimrep}), but it is neither the $\mathbb Z_2$
phase nor a sum of two copies of the $\mathbb Z_1$ phase. It violates the
duality $M_{a^*}=M_a^{\mathsf T}$ of a module category, as $\pi$ is self-dual.

Elimination plan: formalize the duality condition on representations; bound
the number of blocks of an indecomposable representation of each ring
satisfying it; and enumerate the finitely many candidates, as is done for
the Fibonacci ring on two blocks.
\bibliographystyle{alpha}
\bibliography{references}
\end{document}
